% Options for packages loaded elsewhere
\PassOptionsToPackage{unicode}{hyperref}
\PassOptionsToPackage{hyphens}{url}
\PassOptionsToPackage{dvipsnames,svgnames,x11names}{xcolor}
%
\documentclass[
  letterpaper,
  DIV=11,
  numbers=noendperiod]{scrartcl}

\usepackage{amsmath,amssymb}
\usepackage{iftex}
\ifPDFTeX
  \usepackage[T1]{fontenc}
  \usepackage[utf8]{inputenc}
  \usepackage{textcomp} % provide euro and other symbols
\else % if luatex or xetex
  \usepackage{unicode-math}
  \defaultfontfeatures{Scale=MatchLowercase}
  \defaultfontfeatures[\rmfamily]{Ligatures=TeX,Scale=1}
\fi
\usepackage{lmodern}
\ifPDFTeX\else  
    % xetex/luatex font selection
\fi
% Use upquote if available, for straight quotes in verbatim environments
\IfFileExists{upquote.sty}{\usepackage{upquote}}{}
\IfFileExists{microtype.sty}{% use microtype if available
  \usepackage[]{microtype}
  \UseMicrotypeSet[protrusion]{basicmath} % disable protrusion for tt fonts
}{}
\makeatletter
\@ifundefined{KOMAClassName}{% if non-KOMA class
  \IfFileExists{parskip.sty}{%
    \usepackage{parskip}
  }{% else
    \setlength{\parindent}{0pt}
    \setlength{\parskip}{6pt plus 2pt minus 1pt}}
}{% if KOMA class
  \KOMAoptions{parskip=half}}
\makeatother
\usepackage{xcolor}
\setlength{\emergencystretch}{3em} % prevent overfull lines
\setcounter{secnumdepth}{-\maxdimen} % remove section numbering
% Make \paragraph and \subparagraph free-standing
\makeatletter
\ifx\paragraph\undefined\else
  \let\oldparagraph\paragraph
  \renewcommand{\paragraph}{
    \@ifstar
      \xxxParagraphStar
      \xxxParagraphNoStar
  }
  \newcommand{\xxxParagraphStar}[1]{\oldparagraph*{#1}\mbox{}}
  \newcommand{\xxxParagraphNoStar}[1]{\oldparagraph{#1}\mbox{}}
\fi
\ifx\subparagraph\undefined\else
  \let\oldsubparagraph\subparagraph
  \renewcommand{\subparagraph}{
    \@ifstar
      \xxxSubParagraphStar
      \xxxSubParagraphNoStar
  }
  \newcommand{\xxxSubParagraphStar}[1]{\oldsubparagraph*{#1}\mbox{}}
  \newcommand{\xxxSubParagraphNoStar}[1]{\oldsubparagraph{#1}\mbox{}}
\fi
\makeatother


\providecommand{\tightlist}{%
  \setlength{\itemsep}{0pt}\setlength{\parskip}{0pt}}\usepackage{longtable,booktabs,array}
\usepackage{calc} % for calculating minipage widths
% Correct order of tables after \paragraph or \subparagraph
\usepackage{etoolbox}
\makeatletter
\patchcmd\longtable{\par}{\if@noskipsec\mbox{}\fi\par}{}{}
\makeatother
% Allow footnotes in longtable head/foot
\IfFileExists{footnotehyper.sty}{\usepackage{footnotehyper}}{\usepackage{footnote}}
\makesavenoteenv{longtable}
\usepackage{graphicx}
\makeatletter
\def\maxwidth{\ifdim\Gin@nat@width>\linewidth\linewidth\else\Gin@nat@width\fi}
\def\maxheight{\ifdim\Gin@nat@height>\textheight\textheight\else\Gin@nat@height\fi}
\makeatother
% Scale images if necessary, so that they will not overflow the page
% margins by default, and it is still possible to overwrite the defaults
% using explicit options in \includegraphics[width, height, ...]{}
\setkeys{Gin}{width=\maxwidth,height=\maxheight,keepaspectratio}
% Set default figure placement to htbp
\makeatletter
\def\fps@figure{htbp}
\makeatother

\KOMAoption{captions}{tableheading}
\makeatletter
\@ifpackageloaded{caption}{}{\usepackage{caption}}
\AtBeginDocument{%
\ifdefined\contentsname
  \renewcommand*\contentsname{Table of contents}
\else
  \newcommand\contentsname{Table of contents}
\fi
\ifdefined\listfigurename
  \renewcommand*\listfigurename{List of Figures}
\else
  \newcommand\listfigurename{List of Figures}
\fi
\ifdefined\listtablename
  \renewcommand*\listtablename{List of Tables}
\else
  \newcommand\listtablename{List of Tables}
\fi
\ifdefined\figurename
  \renewcommand*\figurename{Figure}
\else
  \newcommand\figurename{Figure}
\fi
\ifdefined\tablename
  \renewcommand*\tablename{Table}
\else
  \newcommand\tablename{Table}
\fi
}
\@ifpackageloaded{float}{}{\usepackage{float}}
\floatstyle{ruled}
\@ifundefined{c@chapter}{\newfloat{codelisting}{h}{lop}}{\newfloat{codelisting}{h}{lop}[chapter]}
\floatname{codelisting}{Listing}
\newcommand*\listoflistings{\listof{codelisting}{List of Listings}}
\makeatother
\makeatletter
\makeatother
\makeatletter
\@ifpackageloaded{caption}{}{\usepackage{caption}}
\@ifpackageloaded{subcaption}{}{\usepackage{subcaption}}
\makeatother

\ifLuaTeX
  \usepackage{selnolig}  % disable illegal ligatures
\fi
\usepackage{bookmark}

\IfFileExists{xurl.sty}{\usepackage{xurl}}{} % add URL line breaks if available
\urlstyle{same} % disable monospaced font for URLs
\hypersetup{
  pdftitle={Ejercicios \textbackslash mathbb\{R\}\^{}2},
  colorlinks=true,
  linkcolor={blue},
  filecolor={Maroon},
  citecolor={Blue},
  urlcolor={Blue},
  pdfcreator={LaTeX via pandoc}}


\title{Ejercicios \(\mathbb{R}^2\)}
\author{}
\date{}

\begin{document}
\maketitle


\section{\texorpdfstring{Taller de Vectores en (\mathbb{R}\^{}2) y
Matrices (2
\times 2)}{Taller de Vectores en (\^{}2) y Matrices (2 )}}\label{taller-de-vectores-en-2-y-matrices-2}

\subsubsection{Ejercicio 1}\label{ejercicio-1}

Sean los vectores (\mathbf{u} = (3, -1)) y (\mathbf{v} = (2, 4)).\\
1. Calcula (\mathbf{u} + \mathbf{v}).\\
2. Calcula (2\mathbf{u} - 3\mathbf{v}).\\
3. Verifica la propiedad distributiva:\\
{[} \alpha(\mathbf{u} + \mathbf{v}) \stackrel{?}{=} \alpha\mathbf{u} +
\alpha\mathbf{v} {]}\\
tomando (\alpha = -1).

\begin{quote}
\textbf{Pista}: Solo necesitas comparar ambos lados de la igualdad para
concluir.
\end{quote}

\subsubsection{Ejercicio 2}\label{ejercicio-2}

Con los mismos vectores (\mathbf{u} = (3, -1)) y (\mathbf{v} = (2,
4)):\\
1. Halla (\mathbf{u} \cdot \mathbf{v}).\\
2. Verifica la propiedad de conmutatividad (\mathbf{u} \cdot \mathbf{v}
= \mathbf{v} \cdot \mathbf{u}).\\
3. \textbf{(Nivel mayor)} Calcula el ángulo (\theta) entre (\mathbf{u})
y (\mathbf{v}) usando la fórmula:\\
{[} \mathbf{u} \cdot \mathbf{v} =
\textbar{}\mathbf{u}\textbar,\textbar{}\mathbf{v}\textbar{}\cos(\theta).
{]}

\subsubsection{Ejercicio 3}\label{ejercicio-3}

Dada la recta que pasa por el punto (\mathbf{P} = (1, -2)) y que
\textbf{tiene como dirección} el vector (\mathbf{d} = (2, 3)):

\begin{enumerate}
\def\labelenumi{\arabic{enumi}.}
\item
  Escribe su \textbf{ecuación paramétrica} en (\mathbb{R}\^{}2).\\
  \textgreater{} Recuerda que la forma paramétrica se describe como:\\
  \textgreater{} {[} \textgreater{} \mathbf{r}(t) = \mathbf{P} +
  t,\mathbf{d}, \quad t \in \mathbb{R}. \textgreater{} {]}
\item
  Verifica que el punto (\mathbf{Q} = (3,1)) \textbf{no} pertenece a
  dicha recta (sustituyendo en la forma paramétrica y comprobando si
  existe un valor de (t) que lo satisfaga).
\item
  Escribe la ecuación de forma cartesiana, es decir de la forma (ax + by
  = c).
\end{enumerate}

\subsubsection{Ejercicio 4}\label{ejercicio-4}

Considera la recta (r\_1) en forma paramétrica: {[} r\_1:
\quad \mathbf{r}(t) = (2, 1) + t , (3, -1). {]}

\begin{enumerate}
\def\labelenumi{\arabic{enumi}.}
\item
  Encuentra la \textbf{ecuación paramétrica} de la recta (r\_2)
  \textbf{paralela} a (r\_1) y que pase por el punto ((4,2)).\\
  \textgreater{} Dos rectas son paralelas si sus vectores direccionales
  son múltiplos.
\item
  Encuentra la \textbf{ecuación paramétrica} de la recta (r\_3)
  \textbf{perpendicular} a (r\_1) que pase por ((4,2)).\\
  \textgreater{} Dos rectas en (\mathbb{R}\^{}2) son perpendiculares si
  sus vectores direccionales (\mathbf{d_1}) y (\mathbf{d_2}) cumplen
  (\mathbf{d_1}\cdot \mathbf{d_2} = 0).
\item
  Determina el \textbf{punto de intersección} entre (r\_3) y la recta
  (r\_1).
\end{enumerate}

\subsubsection{Ejercicio 5}\label{ejercicio-5}

Sea (\{\mathbf{w}\_1, \mathbf{w}\_2\}) con (\mathbf{w}\_1 = (1, 3)) y
(\mathbf{w}\_2 = (2, 6)).

\begin{enumerate}
\def\labelenumi{\arabic{enumi}.}
\tightlist
\item
  Determina si (\mathbf{w}\_1) y (\mathbf{w}\_2) son linealmente
  \textbf{dependientes} o \textbf{independientes}. Explica tu proceso.\\
\item
  Si son dependientes, describe la relación de dependencia. Si son
  independientes, justifica por qué.
\end{enumerate}

\subsubsection{Ejercicio 6}\label{ejercicio-6}

\begin{enumerate}
\def\labelenumi{\arabic{enumi}.}
\tightlist
\item
  Explica en tus propias palabras qué significa que un conjunto de
  vectores sea \textbf{base} de (\mathbb{R}\^{}2).\\
\item
  Dados los vectores (\mathbf{a} = (1, 2)) y (\mathbf{b} = (2, 1)),
  \textbf{¿son una base de} (\mathbb{R}\^{}2)?

  \begin{itemize}
  \tightlist
  \item
    Justifica tu respuesta usando \textbf{independencia lineal}
  \end{itemize}
\end{enumerate}

\subsubsection{Ejercicio 7}\label{ejercicio-7}

\begin{enumerate}
\def\labelenumi{\arabic{enumi}.}
\tightlist
\item
  ¿Por qué (\{(0,0)\}) (el conjunto que solo contiene al vector cero) es
  un \textbf{subespacio} de (\mathbb{R}\^{}2)?\\
\item
  Si (S) es el conjunto de todos los vectores ((x, y)) tales que (x + y
  = 0), demuestra que \textbf{(S) es un subespacio} de
  (\mathbb{R}\^{}2).\\
\item
  Determina una \textbf{base} para el subespacio (S).
\end{enumerate}




\end{document}
